\begin{table*}[t]
\centering
\caption{不同 Rician 噪声水平下各方法的前景区域平均 PSNR/SSIM。所有方法使用完全相同的测试切片、噪声 realization、归一化和前景掩模。}
\label{tab:fair_method_comparison}
\resizebox{\textwidth}{!}{%
\begin{tabular}{lcccccccccccc}
\toprule
Method & \multicolumn{2}{c}{${\sigma}=0.05$} & \multicolumn{2}{c}{${\sigma}=0.10$} & \multicolumn{2}{c}{${\sigma}=0.15$} & \multicolumn{2}{c}{${\sigma}=0.20$} & \multicolumn{2}{c}{${\sigma}=0.25$} & \multicolumn{2}{c}{${\sigma}=0.30$} \\
 & PSNR (dB) & SSIM & PSNR (dB) & SSIM & PSNR (dB) & SSIM & PSNR (dB) & SSIM & PSNR (dB) & SSIM & PSNR (dB) & SSIM \\
\midrule
Noisy & 26.09 & 0.7805 & 20.25 & 0.5800 & 16.99 & 0.4495 & 14.76 & 0.3571 & 13.09 & 0.2876 & 11.81 & 0.2334 \\
DnCNN & 33.07 & 0.9601 & 29.22 & 0.9125 & 26.66 & 0.8594 & 24.58 & 0.8027 & 22.73 & 0.7419 & 20.85 & 0.6717 \\
RicianNet & 30.71 & 0.9414 & 28.02 & 0.8900 & 26.10 & 0.8379 & 24.39 & 0.7842 & 22.66 & 0.7253 & 20.58 & 0.6532 \\
Attention U-Net & 31.41 & 0.9459 & 27.81 & 0.8902 & 25.57 & 0.8338 & 23.82 & 0.7773 & 22.10 & 0.7155 & 20.07 & 0.6393 \\
\bottomrule
\end{tabular}%
}
\end{table*}
